\chapter{Introduction}
\label{chap:introduction}

Large language models (LLMs) such as those based on the transformer
architecture of Vaswani et al.~\cite{Vaswani2017Attention} have transformed the
way people interact with textual information, making it possible to ask
questions in natural language and receive fluent, contextual answers. These
models, however, share a fundamental limitation: the knowledge they rely on is
encoded in their parameters during training and is therefore frozen in time,
limited in scope to the pre-training corpus, and impossible to audit at the
level of individual facts. When confronted with a domain the model has not
seen, or with recent information, an LLM will frequently \emph{hallucinate} ---
producing plausible but factually incorrect statements that are hard to
distinguish from correct ones, especially for non-expert users. This is a
serious obstacle in domains where the cost of a wrong answer is high, such as
legal advice, medical information, or enterprise document search.

\emph{Retrieval-Augmented Generation} (RAG), proposed by Lewis et
al.~\cite{Lewis2020RAG} at NeurIPS 2020, addresses this limitation by
decoupling the \emph{retrieval} of relevant knowledge from the \emph{generation}
of the final answer. Instead of relying exclusively on the model's internal
parameters, a RAG system first searches a dedicated knowledge store for
passages relevant to the user's question, and then feeds these passages to the
language model as additional context. The model is thus grounded in external,
up-to-date, and auditable evidence, and can cite the sources of its claims.
RAG has become the de facto architectural pattern for building domain-specific
question-answering systems on top of general-purpose LLMs.

The retrieval component of a RAG system is typically implemented with a
\emph{vector database}, a specialised data store that indexes high-dimensional
dense embeddings produced by a neural encoder. Given a query embedding, the
database returns the $k$ nearest neighbours under a similarity metric (cosine
or inner product), enabling semantic search that transcends the lexical
matching of traditional full-text search engines. The rapid development of
vector databases in recent years --- from the pioneering GPU-based search
library FAISS~\cite{Johnson2021FAISS} to user-friendly systems such as
ChromaDB, Pinecone, Weaviate, and Qdrant --- has made semantic retrieval
accessible to a much broader audience of developers.

\section{Motivation}
\label{sec:motivation}

Modern organisations accumulate large amounts of unstructured textual
information --- contracts, regulations, menus, nutrition tables, internal
policies, meeting notes --- that is traditionally accessed either by keyword
search or by manually browsing folder hierarchies. Neither approach scales
well: keyword search misses paraphrases and rewordings, while manual browsing
depends on the user knowing where to look. At the same time, general-purpose
chatbots built on LLMs cannot answer questions about private documents,
because they have never seen them.

A RAG system bridges this gap. It ingests the organisation's documents, stores
them in a vector database, and allows any user to ask questions in natural
language, receiving grounded answers with explicit source citations. The
architecture is general enough to serve a law firm looking up clauses in
collective contracts, a restaurant checking the ingredients of its dishes, or
a research group searching its collection of papers.

From the point of view of a student in a Databases master's programme,
building a practical RAG system is an excellent occasion to bring together
several lines of work:
\begin{itemize}
  \item \textbf{Database systems}, in both their relational and vector
        flavours, and the trade-offs between exact lexical retrieval and
        approximate semantic retrieval.
  \item \textbf{Information retrieval}, and in particular the Dense Passage
        Retrieval approach of Karpukhin et al.~\cite{Karpukhin2020DPR} that
        established dense embeddings as a serious alternative to BM25.
  \item \textbf{Natural language processing}, through sentence-level
        embeddings such as those of Reimers and
        Gurevych~\cite{Reimers2019SentenceBERT} and through large generative
        models.
  \item \textbf{Software engineering}, in the form of a layered, well-tested,
        containerised backend and a modern frontend.
\end{itemize}

\section{Objectives}
\label{sec:objectives}

The main objectives of this thesis are:

\begin{enumerate}
  \item To review the state of the art in information retrieval based on
        vector databases, with an emphasis on the RAG paradigm and its main
        building blocks: dense embeddings, approximate nearest-neighbour
        search, and prompt-based grounding of language models.
  \item To design and implement a RAG system that allows authenticated users
        to upload PDF documents, asks questions in natural language against
        their own private document collection, and receive answers with
        verifiable source citations down to the page level.
  \item To enforce strict data isolation between users, so that documents
        uploaded by one account are invisible to every other account at both
        the database and the retrieval layer.
  \item To evaluate the system in terms of retrieval quality, answer
        groundedness, and response latency, and to compare the semantic
        retrieval component against PostgreSQL full-text search on the same
        corpus.
  \item To produce a reference implementation that is clean, containerised,
        test-covered, and reproducible, so that it can serve both as the
        defence of the thesis and as a starting point for future work on
        retrieval-augmented systems.
\end{enumerate}

\section{Main Contributions}
\label{sec:contributions}

The concrete contributions of this work are the following:

\begin{itemize}
  \item A layered, typed, and test-driven FastAPI backend that separates
        concerns between API routes, services, repositories, and domain
        models, providing a blueprint for reproducible RAG applications.
  \item An end-to-end ingestion pipeline that extracts text from PDFs with
        PyMuPDF, chunks it into overlapping segments of approximately 500
        tokens, computes embeddings using Google Gemini, and stores the
        resulting vectors in a ChromaDB collection dedicated to the
        uploading user.
  \item A query pipeline that retrieves the top-$k$ chunks from the user's
        collection, builds a prompt that binds the language model to the
        retrieved evidence, and returns an answer together with
        machine-readable citations $(\text{filename},\ \text{page})$.
  \item A JWT-based authentication layer backed by bcrypt password hashing
        and PostgreSQL persistence, including a set of demo accounts that
        are seeded idempotently on first startup.
  \item A React frontend built with Vite, Tailwind CSS, and the shadcn/ui
        component system, supporting dark mode, document upload, and an
        interactive chat page that displays source badges.
  \item A Docker Compose configuration that orchestrates the entire stack
        (PostgreSQL, FastAPI backend, React frontend served through nginx)
        and exposes it on a single command.
  \item An experimental evaluation comparing semantic retrieval with
        PostgreSQL full-text search on representative document corpora,
        together with measurements of ingestion and query latency.
\end{itemize}

\section{Structure of the Thesis}
\label{sec:structure}

The remainder of this thesis is organised as follows.

\textbf{Chapter~\ref{chap:background}} introduces the theoretical background
necessary to understand the RAG paradigm. It contrasts relational and vector
databases, surveys the evolution of text embeddings from TF-IDF through
Sentence-BERT, and presents the original RAG architecture of Lewis et al.,
placing the present work in the context of related systems.

\textbf{Chapter~\ref{chap:architecture}} describes the high-level architecture
of the proposed system. It details the ingestion and query pipelines, the
per-user isolation strategy at the vector store level, and the authentication
and session management layers.

\textbf{Chapter~\ref{chap:implementation}} dives into the concrete
implementation. It documents the technology stack, the layered backend
organisation, the integration with ChromaDB and Google Gemini, the frontend
structure, and the containerised deployment model.

\textbf{Chapter~\ref{chap:evaluation}} presents the experimental evaluation.
It defines the metrics, describes the evaluation corpora, reports quantitative
results on retrieval quality and latency, and compares the system against a
PostgreSQL full-text-search baseline.

\textbf{Chapter~\ref{chap:conclusions}} summarises the contributions of the
thesis, discusses the limitations of the present implementation, and outlines
directions for future work.
